% =============================================================================
% Proposal_AIWTPGap.tex
% Research proposal: AI Information, Hypothetical Bias & WTP Measurement
% 15-min presentation, ~20 slides
%
% Source spec: quality_reports/specs/research_spec_aiwtp_gap.md (v6)
% Compile: TEXINPUTS=../Preambles:$TEXINPUTS xelatex -interaction=nonstopmode Proposal_AIWTPGap.tex
% =============================================================================

\documentclass[aspectratio=169]{beamer}
% =============================================================================
% Preambles/header.tex — shared LaTeX/Beamer preamble for 03_AIWTPGap
%
% Korean experimental-economics proposal. Compile with XeLaTeX:
%     % =============================================================================
% Preambles/header.tex — shared LaTeX/Beamer preamble for 03_AIWTPGap
%
% Korean experimental-economics proposal. Compile with XeLaTeX:
%     % =============================================================================
% Preambles/header.tex — shared LaTeX/Beamer preamble for 03_AIWTPGap
%
% Korean experimental-economics proposal. Compile with XeLaTeX:
%     \input{header}
% with TEXINPUTS=../Preambles:$TEXINPUTS (the /compile-latex skill does this).
%
% Required font: KoPubWorld Dotum (Medium / Bold). Install from
% https://www.kopus.or.kr/biz/electronic/font.do (free for academic use).
% Fallback options documented under "Korean font setup" below.
% =============================================================================

% -----------------------------------------------------------------------------
% Engine detection + encoding
%
% XeLaTeX is required (Korean rendering via xeCJK depends on it). Under pdfTeX
% xeCJK is unavailable, so we fall back to inputenc and emit a clear warning.
% -----------------------------------------------------------------------------
\usepackage{iftex}
\ifPDFTeX
  \usepackage[utf8]{inputenc}
  \PackageWarningNoLine{header}{%
    pdfTeX detected --- Korean text will not render. Compile with XeLaTeX.}
\fi

\usepackage{amsmath, amssymb, amsthm}
\usepackage{booktabs}
\usepackage{graphicx}

% -----------------------------------------------------------------------------
% Korean font setup (XeLaTeX only)
%
% Loads xeCJK and KoPubWorld Dotum. The fonts live under ~/Library/Fonts/ on
% macOS. We pass file names (not family names) because the family name is
% registered as Korean ("KoPubWorld돋움체") and fontspec is more robust with
% explicit file references.
%
% Fallback ladder if KoPubWorld Dotum is unavailable: comment out the block
% below and uncomment ONE of the alternatives.
% -----------------------------------------------------------------------------
\ifXeTeX
  \usepackage{xeCJK}
  \xeCJKsetup{%
    CJKspace=true,           % automatic CJK-Latin spacing
    AutoFakeSlant=0.18,      % italic-style fallback when no italic cut exists
    CJKecglue={\hskip 0.18em plus 0.04em minus 0.04em}%
  }

  % --- KoPubWorld Dotum (primary) -------------------------------------------
  % macOS Font Book registers these by PostScript name (no spaces). XeLaTeX's
  % fontspec finds them via the system font registry without needing a
  % Path= override. If your machine doesn't find them, install KoPubWorld
  % Dotum (Medium / Bold / Light) into ~/Library/Fonts/ and rebuild the
  % font cache, or comment this block out and uncomment the Apple SD Gothic
  % Neo fallback below.
  \setCJKmainfont[%
    BoldFont=KoPubWorldDotumBold,
    AutoFakeSlant=0.18
  ]{KoPubWorldDotumMedium}
  \setCJKsansfont[%
    BoldFont=KoPubWorldDotumBold,
    AutoFakeSlant=0.18
  ]{KoPubWorldDotumMedium}
  \setCJKmonofont{KoPubWorldDotumLight}

  % --- Apple SD Gothic Neo fallback (uncomment if KoPub unavailable) --------
  % \setCJKmainfont{Apple SD Gothic Neo}
  % \setCJKsansfont{Apple SD Gothic Neo}
  % \setCJKmonofont{Apple SD Gothic Neo}

  % --- Pretendard fallback (uncomment if installed) -------------------------
  % \setCJKmainfont{Pretendard}
  % \setCJKsansfont{Pretendard}
  % \setCJKmonofont{Pretendard}

  % Slightly looser line spread improves Hangul readability on slides
  \linespread{1.15}
\fi

% -----------------------------------------------------------------------------
% PALETTE — Beamer theme + TikZ-snippet colors.
% Load xcolor and define colors BEFORE anything that references them
% (notably hyperref, hypersetup, and the Beamer color assignments).
% -----------------------------------------------------------------------------
\usepackage{xcolor}

\definecolor{primary-blue}{HTML}{012169}      % headings, accents
\definecolor{primary-gold}{HTML}{B9975B}      % emphasis, borders
\definecolor{highlight-yellow}{HTML}{F2A900}  % markers, alerts
\definecolor{light-bg}{HTML}{E8EDF5}          % title / section backgrounds
\definecolor{jet}{HTML}{1A1A1A}               % body text

% Semantic colors (used in diagrams and annotations)
\definecolor{positive}{HTML}{15803D}          % good / observed / identified
\definecolor{negative}{HTML}{B91C1C}          % bad / problematic / bias
\definecolor{neutral}{HTML}{525252}           % reference / context

% Highlight accents (match .hi-* classes in the SCSS)
\definecolor{hi-slate}{HTML}{314F4F}
\definecolor{hi-green}{HTML}{2E8B57}
\definecolor{hi-red}{HTML}{C41E3A}

% Semantic aliases so skill templates and snippets can write
% \textcolor{accent}{...} without hard-coding "primary-blue".
\colorlet{accent}{primary-blue}
\colorlet{accent2}{primary-gold}

% -----------------------------------------------------------------------------
% Hyperref — loaded AFTER xcolor + palette so link colors resolve.
% -----------------------------------------------------------------------------
\usepackage{hyperref}
\hypersetup{colorlinks=true, linkcolor=primary-blue, urlcolor=primary-blue,
            citecolor=primary-blue, pdfborder={0 0 0}}

% -----------------------------------------------------------------------------
% Beamer theme assignments (only apply under Beamer).
%
% We detect Beamer via \@ifclassloaded{beamer}, which is the documented,
% non-internal way. This must be wrapped in \makeatletter/\makeatother
% because \@ifclassloaded contains an @-catcode character.
% -----------------------------------------------------------------------------
\makeatletter
\@ifclassloaded{beamer}{%
  \usefonttheme{professionalfonts}
  \setbeamercolor{structure}{fg=primary-blue}
  \setbeamercolor{frametitle}{fg=primary-blue}
  \setbeamercolor{title}{fg=primary-blue}
  \setbeamercolor{subtitle}{fg=primary-gold}
  \setbeamercolor{author}{fg=primary-blue}
  \setbeamercolor{institute}{fg=neutral}
  \setbeamercolor{date}{fg=primary-gold}
  \setbeamercolor{itemize item}{fg=primary-blue}
  \setbeamercolor{itemize subitem}{fg=primary-gold}
  \setbeamercolor{alerted text}{fg=primary-gold}
  \setbeamercolor{block title}{fg=white, bg=primary-blue}
  \setbeamercolor{block body}{bg=light-bg}

  % Minimal footer: page X / N only, no navigation clutter
  \setbeamertemplate{navigation symbols}{}
  \setbeamertemplate{footline}{%
    \hfill{\color{neutral}\footnotesize\insertframenumber/\inserttotalframenumber}\hspace{0.7em}\vspace{0.3em}}
}{}
\makeatother

% -----------------------------------------------------------------------------
% TikZ setup — libraries the snippet gallery uses
% -----------------------------------------------------------------------------
\usepackage{tikz}
\usetikzlibrary{arrows.meta, positioning, calc, decorations.pathreplacing,
                fit, shapes.geometric, backgrounds}

% Shared TikZ styles — snippets and hand-written diagrams can pick these up
% via `\begin{tikzpicture}[every node/.style={...}]` or by naming specific
% styles (e.g., `\node[dag-node] (X) at (0,0) {X};`).
\tikzset{
  % Node styles for causal DAGs and similar
  dag-node/.style = {
    draw=primary-blue, fill=light-bg, rounded corners=2pt,
    minimum width=1.2cm, minimum height=0.9cm, align=center, font=\small
  },
  decision-node/.style = {
    draw=primary-gold, fill=white, diamond, aspect=1.8,
    minimum width=1.8cm, minimum height=1.2cm, align=center, font=\small,
    inner sep=1pt
  },
  % Edge styles — observed vs counterfactual
  observed-edge/.style = {-{Latex[length=2.5mm]}, thick, draw=jet},
  counterfactual-edge/.style = {-{Latex[length=2.5mm]}, thick, draw=neutral, dashed},
  confound-edge/.style = {-{Latex[length=2.5mm]}, thick, draw=negative, dashed},
  % Data-point styles for plots
  observed-dot/.style = {fill=primary-blue, draw=primary-blue, circle, inner sep=1.2pt},
  counterfactual-dot/.style = {fill=white, draw=neutral, circle, inner sep=1.2pt}
}

% -----------------------------------------------------------------------------
% Convenience macros
% -----------------------------------------------------------------------------
\newcommand{\muted}[1]{\textcolor{neutral}{#1}}
\newcommand{\key}[1]{\textcolor{primary-gold}{\textbf{#1}}}
\newcommand{\good}[1]{\textcolor{positive}{#1}}
\newcommand{\bad}[1]{\textcolor{negative}{#1}}

% Standout frame macro: full-bleed dark background for section transitions.
% Expand: \transitionslide{Your transition title}
\newcommand{\transitionslide}[1]{%
  \begingroup
  \setbeamercolor{background canvas}{bg=primary-blue}
  \begin{frame}[plain]
    \centering
    \vfill
    {\usebeamerfont{title}\color{white}\Large #1\par}
    \vfill
  \end{frame}
  \endgroup
}

% with TEXINPUTS=../Preambles:$TEXINPUTS (the /compile-latex skill does this).
%
% Required font: KoPubWorld Dotum (Medium / Bold). Install from
% https://www.kopus.or.kr/biz/electronic/font.do (free for academic use).
% Fallback options documented under "Korean font setup" below.
% =============================================================================

% -----------------------------------------------------------------------------
% Engine detection + encoding
%
% XeLaTeX is required (Korean rendering via xeCJK depends on it). Under pdfTeX
% xeCJK is unavailable, so we fall back to inputenc and emit a clear warning.
% -----------------------------------------------------------------------------
\usepackage{iftex}
\ifPDFTeX
  \usepackage[utf8]{inputenc}
  \PackageWarningNoLine{header}{%
    pdfTeX detected --- Korean text will not render. Compile with XeLaTeX.}
\fi

\usepackage{amsmath, amssymb, amsthm}
\usepackage{booktabs}
\usepackage{graphicx}

% -----------------------------------------------------------------------------
% Korean font setup (XeLaTeX only)
%
% Loads xeCJK and KoPubWorld Dotum. The fonts live under ~/Library/Fonts/ on
% macOS. We pass file names (not family names) because the family name is
% registered as Korean ("KoPubWorld돋움체") and fontspec is more robust with
% explicit file references.
%
% Fallback ladder if KoPubWorld Dotum is unavailable: comment out the block
% below and uncomment ONE of the alternatives.
% -----------------------------------------------------------------------------
\ifXeTeX
  \usepackage{xeCJK}
  \xeCJKsetup{%
    CJKspace=true,           % automatic CJK-Latin spacing
    AutoFakeSlant=0.18,      % italic-style fallback when no italic cut exists
    CJKecglue={\hskip 0.18em plus 0.04em minus 0.04em}%
  }

  % --- KoPubWorld Dotum (primary) -------------------------------------------
  % macOS Font Book registers these by PostScript name (no spaces). XeLaTeX's
  % fontspec finds them via the system font registry without needing a
  % Path= override. If your machine doesn't find them, install KoPubWorld
  % Dotum (Medium / Bold / Light) into ~/Library/Fonts/ and rebuild the
  % font cache, or comment this block out and uncomment the Apple SD Gothic
  % Neo fallback below.
  \setCJKmainfont[%
    BoldFont=KoPubWorldDotumBold,
    AutoFakeSlant=0.18
  ]{KoPubWorldDotumMedium}
  \setCJKsansfont[%
    BoldFont=KoPubWorldDotumBold,
    AutoFakeSlant=0.18
  ]{KoPubWorldDotumMedium}
  \setCJKmonofont{KoPubWorldDotumLight}

  % --- Apple SD Gothic Neo fallback (uncomment if KoPub unavailable) --------
  % \setCJKmainfont{Apple SD Gothic Neo}
  % \setCJKsansfont{Apple SD Gothic Neo}
  % \setCJKmonofont{Apple SD Gothic Neo}

  % --- Pretendard fallback (uncomment if installed) -------------------------
  % \setCJKmainfont{Pretendard}
  % \setCJKsansfont{Pretendard}
  % \setCJKmonofont{Pretendard}

  % Slightly looser line spread improves Hangul readability on slides
  \linespread{1.15}
\fi

% -----------------------------------------------------------------------------
% PALETTE — Beamer theme + TikZ-snippet colors.
% Load xcolor and define colors BEFORE anything that references them
% (notably hyperref, hypersetup, and the Beamer color assignments).
% -----------------------------------------------------------------------------
\usepackage{xcolor}

\definecolor{primary-blue}{HTML}{012169}      % headings, accents
\definecolor{primary-gold}{HTML}{B9975B}      % emphasis, borders
\definecolor{highlight-yellow}{HTML}{F2A900}  % markers, alerts
\definecolor{light-bg}{HTML}{E8EDF5}          % title / section backgrounds
\definecolor{jet}{HTML}{1A1A1A}               % body text

% Semantic colors (used in diagrams and annotations)
\definecolor{positive}{HTML}{15803D}          % good / observed / identified
\definecolor{negative}{HTML}{B91C1C}          % bad / problematic / bias
\definecolor{neutral}{HTML}{525252}           % reference / context

% Highlight accents (match .hi-* classes in the SCSS)
\definecolor{hi-slate}{HTML}{314F4F}
\definecolor{hi-green}{HTML}{2E8B57}
\definecolor{hi-red}{HTML}{C41E3A}

% Semantic aliases so skill templates and snippets can write
% \textcolor{accent}{...} without hard-coding "primary-blue".
\colorlet{accent}{primary-blue}
\colorlet{accent2}{primary-gold}

% -----------------------------------------------------------------------------
% Hyperref — loaded AFTER xcolor + palette so link colors resolve.
% -----------------------------------------------------------------------------
\usepackage{hyperref}
\hypersetup{colorlinks=true, linkcolor=primary-blue, urlcolor=primary-blue,
            citecolor=primary-blue, pdfborder={0 0 0}}

% -----------------------------------------------------------------------------
% Beamer theme assignments (only apply under Beamer).
%
% We detect Beamer via \@ifclassloaded{beamer}, which is the documented,
% non-internal way. This must be wrapped in \makeatletter/\makeatother
% because \@ifclassloaded contains an @-catcode character.
% -----------------------------------------------------------------------------
\makeatletter
\@ifclassloaded{beamer}{%
  \usefonttheme{professionalfonts}
  \setbeamercolor{structure}{fg=primary-blue}
  \setbeamercolor{frametitle}{fg=primary-blue}
  \setbeamercolor{title}{fg=primary-blue}
  \setbeamercolor{subtitle}{fg=primary-gold}
  \setbeamercolor{author}{fg=primary-blue}
  \setbeamercolor{institute}{fg=neutral}
  \setbeamercolor{date}{fg=primary-gold}
  \setbeamercolor{itemize item}{fg=primary-blue}
  \setbeamercolor{itemize subitem}{fg=primary-gold}
  \setbeamercolor{alerted text}{fg=primary-gold}
  \setbeamercolor{block title}{fg=white, bg=primary-blue}
  \setbeamercolor{block body}{bg=light-bg}

  % Minimal footer: page X / N only, no navigation clutter
  \setbeamertemplate{navigation symbols}{}
  \setbeamertemplate{footline}{%
    \hfill{\color{neutral}\footnotesize\insertframenumber/\inserttotalframenumber}\hspace{0.7em}\vspace{0.3em}}
}{}
\makeatother

% -----------------------------------------------------------------------------
% TikZ setup — libraries the snippet gallery uses
% -----------------------------------------------------------------------------
\usepackage{tikz}
\usetikzlibrary{arrows.meta, positioning, calc, decorations.pathreplacing,
                fit, shapes.geometric, backgrounds}

% Shared TikZ styles — snippets and hand-written diagrams can pick these up
% via `\begin{tikzpicture}[every node/.style={...}]` or by naming specific
% styles (e.g., `\node[dag-node] (X) at (0,0) {X};`).
\tikzset{
  % Node styles for causal DAGs and similar
  dag-node/.style = {
    draw=primary-blue, fill=light-bg, rounded corners=2pt,
    minimum width=1.2cm, minimum height=0.9cm, align=center, font=\small
  },
  decision-node/.style = {
    draw=primary-gold, fill=white, diamond, aspect=1.8,
    minimum width=1.8cm, minimum height=1.2cm, align=center, font=\small,
    inner sep=1pt
  },
  % Edge styles — observed vs counterfactual
  observed-edge/.style = {-{Latex[length=2.5mm]}, thick, draw=jet},
  counterfactual-edge/.style = {-{Latex[length=2.5mm]}, thick, draw=neutral, dashed},
  confound-edge/.style = {-{Latex[length=2.5mm]}, thick, draw=negative, dashed},
  % Data-point styles for plots
  observed-dot/.style = {fill=primary-blue, draw=primary-blue, circle, inner sep=1.2pt},
  counterfactual-dot/.style = {fill=white, draw=neutral, circle, inner sep=1.2pt}
}

% -----------------------------------------------------------------------------
% Convenience macros
% -----------------------------------------------------------------------------
\newcommand{\muted}[1]{\textcolor{neutral}{#1}}
\newcommand{\key}[1]{\textcolor{primary-gold}{\textbf{#1}}}
\newcommand{\good}[1]{\textcolor{positive}{#1}}
\newcommand{\bad}[1]{\textcolor{negative}{#1}}

% Standout frame macro: full-bleed dark background for section transitions.
% Expand: \transitionslide{Your transition title}
\newcommand{\transitionslide}[1]{%
  \begingroup
  \setbeamercolor{background canvas}{bg=primary-blue}
  \begin{frame}[plain]
    \centering
    \vfill
    {\usebeamerfont{title}\color{white}\Large #1\par}
    \vfill
  \end{frame}
  \endgroup
}

% with TEXINPUTS=../Preambles:$TEXINPUTS (the /compile-latex skill does this).
%
% Required font: KoPubWorld Dotum (Medium / Bold). Install from
% https://www.kopus.or.kr/biz/electronic/font.do (free for academic use).
% Fallback options documented under "Korean font setup" below.
% =============================================================================

% -----------------------------------------------------------------------------
% Engine detection + encoding
%
% XeLaTeX is required (Korean rendering via xeCJK depends on it). Under pdfTeX
% xeCJK is unavailable, so we fall back to inputenc and emit a clear warning.
% -----------------------------------------------------------------------------
\usepackage{iftex}
\ifPDFTeX
  \usepackage[utf8]{inputenc}
  \PackageWarningNoLine{header}{%
    pdfTeX detected --- Korean text will not render. Compile with XeLaTeX.}
\fi

\usepackage{amsmath, amssymb, amsthm}
\usepackage{booktabs}
\usepackage{graphicx}

% -----------------------------------------------------------------------------
% Korean font setup (XeLaTeX only)
%
% Loads xeCJK and KoPubWorld Dotum. The fonts live under ~/Library/Fonts/ on
% macOS. We pass file names (not family names) because the family name is
% registered as Korean ("KoPubWorld돋움체") and fontspec is more robust with
% explicit file references.
%
% Fallback ladder if KoPubWorld Dotum is unavailable: comment out the block
% below and uncomment ONE of the alternatives.
% -----------------------------------------------------------------------------
\ifXeTeX
  \usepackage{xeCJK}
  \xeCJKsetup{%
    CJKspace=true,           % automatic CJK-Latin spacing
    AutoFakeSlant=0.18,      % italic-style fallback when no italic cut exists
    CJKecglue={\hskip 0.18em plus 0.04em minus 0.04em}%
  }

  % --- KoPubWorld Dotum (primary) -------------------------------------------
  % macOS Font Book registers these by PostScript name (no spaces). XeLaTeX's
  % fontspec finds them via the system font registry without needing a
  % Path= override. If your machine doesn't find them, install KoPubWorld
  % Dotum (Medium / Bold / Light) into ~/Library/Fonts/ and rebuild the
  % font cache, or comment this block out and uncomment the Apple SD Gothic
  % Neo fallback below.
  \setCJKmainfont[%
    BoldFont=KoPubWorldDotumBold,
    AutoFakeSlant=0.18
  ]{KoPubWorldDotumMedium}
  \setCJKsansfont[%
    BoldFont=KoPubWorldDotumBold,
    AutoFakeSlant=0.18
  ]{KoPubWorldDotumMedium}
  \setCJKmonofont{KoPubWorldDotumLight}

  % --- Apple SD Gothic Neo fallback (uncomment if KoPub unavailable) --------
  % \setCJKmainfont{Apple SD Gothic Neo}
  % \setCJKsansfont{Apple SD Gothic Neo}
  % \setCJKmonofont{Apple SD Gothic Neo}

  % --- Pretendard fallback (uncomment if installed) -------------------------
  % \setCJKmainfont{Pretendard}
  % \setCJKsansfont{Pretendard}
  % \setCJKmonofont{Pretendard}

  % Slightly looser line spread improves Hangul readability on slides
  \linespread{1.15}
\fi

% -----------------------------------------------------------------------------
% PALETTE — Beamer theme + TikZ-snippet colors.
% Load xcolor and define colors BEFORE anything that references them
% (notably hyperref, hypersetup, and the Beamer color assignments).
% -----------------------------------------------------------------------------
\usepackage{xcolor}

\definecolor{primary-blue}{HTML}{012169}      % headings, accents
\definecolor{primary-gold}{HTML}{B9975B}      % emphasis, borders
\definecolor{highlight-yellow}{HTML}{F2A900}  % markers, alerts
\definecolor{light-bg}{HTML}{E8EDF5}          % title / section backgrounds
\definecolor{jet}{HTML}{1A1A1A}               % body text

% Semantic colors (used in diagrams and annotations)
\definecolor{positive}{HTML}{15803D}          % good / observed / identified
\definecolor{negative}{HTML}{B91C1C}          % bad / problematic / bias
\definecolor{neutral}{HTML}{525252}           % reference / context

% Highlight accents (match .hi-* classes in the SCSS)
\definecolor{hi-slate}{HTML}{314F4F}
\definecolor{hi-green}{HTML}{2E8B57}
\definecolor{hi-red}{HTML}{C41E3A}

% Semantic aliases so skill templates and snippets can write
% \textcolor{accent}{...} without hard-coding "primary-blue".
\colorlet{accent}{primary-blue}
\colorlet{accent2}{primary-gold}

% -----------------------------------------------------------------------------
% Hyperref — loaded AFTER xcolor + palette so link colors resolve.
% -----------------------------------------------------------------------------
\usepackage{hyperref}
\hypersetup{colorlinks=true, linkcolor=primary-blue, urlcolor=primary-blue,
            citecolor=primary-blue, pdfborder={0 0 0}}

% -----------------------------------------------------------------------------
% Beamer theme assignments (only apply under Beamer).
%
% We detect Beamer via \@ifclassloaded{beamer}, which is the documented,
% non-internal way. This must be wrapped in \makeatletter/\makeatother
% because \@ifclassloaded contains an @-catcode character.
% -----------------------------------------------------------------------------
\makeatletter
\@ifclassloaded{beamer}{%
  \usefonttheme{professionalfonts}
  \setbeamercolor{structure}{fg=primary-blue}
  \setbeamercolor{frametitle}{fg=primary-blue}
  \setbeamercolor{title}{fg=primary-blue}
  \setbeamercolor{subtitle}{fg=primary-gold}
  \setbeamercolor{author}{fg=primary-blue}
  \setbeamercolor{institute}{fg=neutral}
  \setbeamercolor{date}{fg=primary-gold}
  \setbeamercolor{itemize item}{fg=primary-blue}
  \setbeamercolor{itemize subitem}{fg=primary-gold}
  \setbeamercolor{alerted text}{fg=primary-gold}
  \setbeamercolor{block title}{fg=white, bg=primary-blue}
  \setbeamercolor{block body}{bg=light-bg}

  % Minimal footer: page X / N only, no navigation clutter
  \setbeamertemplate{navigation symbols}{}
  \setbeamertemplate{footline}{%
    \hfill{\color{neutral}\footnotesize\insertframenumber/\inserttotalframenumber}\hspace{0.7em}\vspace{0.3em}}
}{}
\makeatother

% -----------------------------------------------------------------------------
% TikZ setup — libraries the snippet gallery uses
% -----------------------------------------------------------------------------
\usepackage{tikz}
\usetikzlibrary{arrows.meta, positioning, calc, decorations.pathreplacing,
                fit, shapes.geometric, backgrounds}

% Shared TikZ styles — snippets and hand-written diagrams can pick these up
% via `\begin{tikzpicture}[every node/.style={...}]` or by naming specific
% styles (e.g., `\node[dag-node] (X) at (0,0) {X};`).
\tikzset{
  % Node styles for causal DAGs and similar
  dag-node/.style = {
    draw=primary-blue, fill=light-bg, rounded corners=2pt,
    minimum width=1.2cm, minimum height=0.9cm, align=center, font=\small
  },
  decision-node/.style = {
    draw=primary-gold, fill=white, diamond, aspect=1.8,
    minimum width=1.8cm, minimum height=1.2cm, align=center, font=\small,
    inner sep=1pt
  },
  % Edge styles — observed vs counterfactual
  observed-edge/.style = {-{Latex[length=2.5mm]}, thick, draw=jet},
  counterfactual-edge/.style = {-{Latex[length=2.5mm]}, thick, draw=neutral, dashed},
  confound-edge/.style = {-{Latex[length=2.5mm]}, thick, draw=negative, dashed},
  % Data-point styles for plots
  observed-dot/.style = {fill=primary-blue, draw=primary-blue, circle, inner sep=1.2pt},
  counterfactual-dot/.style = {fill=white, draw=neutral, circle, inner sep=1.2pt}
}

% -----------------------------------------------------------------------------
% Convenience macros
% -----------------------------------------------------------------------------
\newcommand{\muted}[1]{\textcolor{neutral}{#1}}
\newcommand{\key}[1]{\textcolor{primary-gold}{\textbf{#1}}}
\newcommand{\good}[1]{\textcolor{positive}{#1}}
\newcommand{\bad}[1]{\textcolor{negative}{#1}}

% Standout frame macro: full-bleed dark background for section transitions.
% Expand: \transitionslide{Your transition title}
\newcommand{\transitionslide}[1]{%
  \begingroup
  \setbeamercolor{background canvas}{bg=primary-blue}
  \begin{frame}[plain]
    \centering
    \vfill
    {\usebeamerfont{title}\color{white}\Large #1\par}
    \vfill
  \end{frame}
  \endgroup
}


% Local override: KoPubWorld Dotum is no longer detected on this system as
% of 2026-05-28. Fall back to Apple SD Gothic Neo (macOS system font, covers
% both Korean and basic Latin glyphs) for all main/sans/CJK contexts.
\setCJKmainfont{Apple SD Gothic Neo}
\setCJKsansfont{Apple SD Gothic Neo}
\setCJKmonofont{Apple SD Gothic Neo}
\setmainfont{Apple SD Gothic Neo}
\setsansfont{Apple SD Gothic Neo}

\usepackage{tikz}
\usetikzlibrary{arrows.meta, positioning, shapes.geometric, fit, calc}

\title{AI 정보 제공과 WTP 측정의 가상편의}
\subtitle{Algorithm Appreciation은 정보의 보정인가, 왜곡인가?}
\author{[저자명]}
\institute{[소속]}
\date{\today}

\begin{document}

% =============================================================================
% Part I: Motivation & Research Question (3 slides)
% =============================================================================

\frame{\titlepage}

% --- Slide 2 ----------------------------------------------------------------
\begin{frame}{왜 이 연구인가 --- 두 가지 사실의 충돌}
  \begin{itemize}
    \item \textbf{사실 1 (Fox et al.\ 1998 CVM-X):}\\
      식품·환경 정책의 WTP 측정에서 가상편의(hypothetical bias)는 robust한 현상.\\
      $\mathrm{WTP}_{\text{CVM}} \approx 1.5\text{--}1.8 \times \mathrm{WTP}_{\text{auction}}$.
    \item \textbf{사실 2 (Logg et al.\ 2019 / Dietvorst et al.\ 2015):}\\
      소비자는 AI 조언에 인간 전문가보다 더 높은 가중치를 부여한다(Algorithm Appreciation, WOA $d \approx 0.4$).\\
      그러나 AI 오류 목격 후 AI 선택률이 65\,\% $\rightarrow$ 23\,\%로 급락한다(Algorithm Aversion).
    \item \textbf{Corrigan et al.\ (2012):} 정보 제공이 demand revelation을 오히려 \emph{악화}시킬 수 있다.
  \end{itemize}
  \vspace{0.4em}
  \key{질문: AI 정보는 WTP 측정의 가상편의를 \emph{줄이는가}, 아니면 \emph{왜곡}시키는가?}
\end{frame}

% --- Slide 3 ----------------------------------------------------------------
\begin{frame}{연구 질문과 핵심 기여}
  \textbf{Research Question:}
  \begin{quote}
    Does AI-personalized product information \good{reduce} hypothetical bias in WTP elicitation,
    or does it \bad{distort} valuations through anchoring, information overload, or algorithm bias?
  \end{quote}

  \textbf{세부 질문:}
  \begin{itemize}
    \item Q1. AI 개인화 정보 vs 일반 정보 (정보원 효과)
    \item Q2. 정보량 증가 효과 (과부하 임계점)
    \item Q3. AI $\times$ 고정보량 상호작용 (앵커링·과부하·자동화 편향)
    \item Q4. 인지능력(CRT)의 조절 효과 (탐색적)
  \end{itemize}

  \vspace{0.3em}
  \key{핵심 기여:} Fox CVM-X 라인의 21세기 AI 재해석 --- AI를 \emph{calibration tool}로 평가한 최초의 lab 실험.
\end{frame}

% =============================================================================
% Part II: Theoretical Framework (5 slides)
% =============================================================================

% --- Slide 4 ----------------------------------------------------------------
\begin{frame}{이론적 배경 (1) --- 가상편의(Hypothetical Bias)}
  \[
    \mathrm{Bias} \;=\; \mathrm{WTP}_{\text{CVM}} \;-\; \mathrm{WTP}_{\text{auction}}
  \]
  \begin{itemize}
    \item CVM 응답은 인센티브 비호환적 $\rightarrow$ 예산·참여 제약 부재 $\rightarrow$ WTP 과장.
    \item 실험경매는 인센티브 호환적이나 고비용·소규모.
    \item Fox et al.\ (1998): 보정계수 0.55--0.69 (실재 근거로만 활용; 분석에서 직접 사용 안 함).
  \end{itemize}

  \vspace{0.3em}
  \muted{본 연구는 endowment effect의 reference-dependent 논쟁(Köszegi--Rabin 2006; Heffetz--List 2014; Campos-Mercade et al.\ 2026)에 참여하지 않는다.
  WTP--WTA 갭이 아닌 \emph{CVM--경매 갭}이 본 연구의 종속변수.}
\end{frame}

% --- Slide 5 ----------------------------------------------------------------
\begin{frame}{이론적 배경 (2) --- 채널 A: Bayesian 업데이트}
  \textbf{예측: 갭 감소}
  \begin{align*}
    \text{AI 정보} \;&\rightarrow\; \text{credence good 효능 불확실성 } \downarrow \\
    &\rightarrow\; \text{경매 WTP } \uparrow \;\text{(진짜 가치 수렴)} \\
    &\rightarrow\; \text{Bias} \;\downarrow
  \end{align*}

  \textbf{비대칭 이동 메커니즘:}
  \begin{itemize}
    \item 경매: commitment cost $> 0$ (Zhao--Kling 2001) $\rightarrow$ 정보 정밀도 $\uparrow$ $\rightarrow$ 입찰 정확도 $\uparrow$.
    \item CVM: commitment cost $\approx 0$ (가상) $\rightarrow$ 정보의 한계 효과 제한.
    \item 예측: $\beta_{\text{auction}} > \beta_{\text{CVM}}$ $\rightarrow$ \textbf{SUR 검정}으로 확인 (Slide 16).
  \end{itemize}
\end{frame}

% --- Slide 6 ----------------------------------------------------------------
\begin{frame}{이론적 배경 (3) --- 채널 B-1, B-2}
  \textbf{B-1. 앵커링 (Anchoring)}
  \begin{itemize}
    \item AI 추정 가치 ``XX,XXX원'' $\rightarrow$ 금전적 앵커 형성.
    \item 입찰이 AI 제시가로 수렴: $\mathrm{AnchorDev} = \mathrm{Bid} - \mathrm{AI\_price\_estimate} \approx 0$.
    \item Bias 방향: $\mathrm{AI\_price} < \text{진짜 WTP}$이면 갭 $\downarrow$, $>$이면 갭 $\uparrow$ (확인 불확실).
  \end{itemize}

  \vspace{0.4em}
  \textbf{B-2. 정보 과부하 (Information Overload)}
  \begin{itemize}
    \item AI 상세 정보 $\rightarrow$ 처리 용량 초과 $\rightarrow$ System 1 (휴리스틱) 작동.
    \item 입찰 \emph{분산} 증가 (mean 이동 없음; overload $\neq$ anchoring).
    \item Lee et al.\ (2020): 저CRT 집단 입찰 분산이 3배.
  \end{itemize}
\end{frame}

% --- Slide 7 ----------------------------------------------------------------
\begin{frame}{이론적 배경 (4) --- 채널 B-3: Algorithm Appreciation/Aversion}
  \textbf{B-3. 자동화 편향}
  \begin{itemize}
    \item \textbf{Logg et al.\ (2019, \emph{OBHDP}):} Appreciation = AI 조언에 인간보다 \emph{더} 가중치 (WOA $d \approx 0.4$).
    \item \textbf{Dietvorst et al.\ (2015, \emph{JEP:General}):} Aversion = AI 오류 목격 후 AI 선택률 65\,\% $\rightarrow$ 23\,\% 급락.
  \end{itemize}

  \vspace{0.4em}
  \key{Logg-Dietvorst Boundary Condition (Logg 2019, p.91 직접 인용):}\\
  Appreciation은 \emph{default} (오류 관측 전); Aversion은 \emph{조건부} (오류 관측 후).\\
  \textbf{$\Rightarrow$ Phase 1 중립화로 default 유지 (Slide 9).}

  \vspace{0.4em}
  \textbf{AA $\rightarrow$ WTP 연결:} AI 정보가 가치 평가의 \emph{프레임}을 설정 $\rightarrow$ WTP가 AI 제시가 방향으로 \emph{무의식적} 이동 (B-1 앵커링과 구분).
\end{frame}

% --- Slide 8 ----------------------------------------------------------------
\begin{frame}{이론적 배경 (5) --- 핵심 예측표}
  \small
  \begin{tabular}{lccc}
    \hline
    \textbf{조건} & \textbf{채널 A (Bayesian)} & \textbf{채널 B (행동경제)} & \textbf{식별} \\
    \hline
    T1: Generic $\times$ Low & 갭 소폭 $\downarrow$ & 변화 미미 & T1 vs T0 \\
    T2: Generic $\times$ High & 갭 $\downarrow$ & 과부하 분산 $\uparrow$ & T2 vs T1 \\
    T3: AI $\times$ Low & 갭 $\downarrow$ & Appreciation 변동 & T3 vs T1 \\
    \textbf{T4: AI $\times$ High} & 갭 더 $\downarrow$ & \textbf{앵커링+과부하+편향} $\rightarrow$ 갭 $\uparrow$? & \textbf{T4 vs T3, T4 vs T2} \\
    \hline
  \end{tabular}

  \vspace{0.6em}
  \key{핵심 조건: T4 (AI $\times$ High)} --- Bayesian과 채널 B의 \emph{어느 쪽이 지배}하는가가 본 연구의 핵심.\\
  결과가 음수든 양수든 \emph{해석 가능}.
\end{frame}

% =============================================================================
% Part III: Experimental Design (6 slides)
% =============================================================================

% --- Slide 9: TikZ 흐름도 ---------------------------------------------------
\begin{frame}{실험 설계 (1) --- 2-Phase 구조}
  % Coordinate map (P2 compliant):
  %   Phase1 box at (0, 3.0)    -- 중립화, 모든 집단 동일
  %   Phase2 box at (0, 0.0)    -- 처치 → CVM → filler → 경매
  %   Outcome  at  (0, -2.5)    -- Bias = WTP_CVM - WTP_auction
  \begin{center}
  \begin{tikzpicture}[
    scale=1.0, every node/.style={scale=1.0},
    box/.style={draw, rounded corners, minimum width=11cm, minimum height=1.4cm,
                text width=10.6cm, align=center, font=\small},
    arrow/.style={-{Stealth[length=3mm]}, thick}
  ]
    \node[box, fill=blue!8] (P1) at (0, 3.0)
      {\textbf{Phase 1: Induced Value (벤치마크 / 중립)}\\
       모든 집단 동일 절차 --- 메커니즘 교육 + 연습 5라운드 --- AI 정보 \emph{없음}};

    \node[box, fill=orange!10] (P2) at (0, 0.0)
      {\textbf{Phase 2: Homegrown Value (본 실험)}\\
       처치 (T0--T4) $\rightarrow$ CVM 응답 $\rightarrow$ Filler (5분) $\rightarrow$ Random nth-price 경매};

    \node[box, fill=green!12] (OUT) at (0, -2.5)
      {\textbf{측정: Bias} $=$ $\mathrm{WTP}_{\text{CVM}} - \mathrm{WTP}_{\text{auction}}$\\
       + AnchorDev $=$ Bid $-$ AI\_price\_estimate (T3/T4)};

    \draw[arrow] (P1) -- node[right, font=\footnotesize] {동일 피험자} (P2);
    \draw[arrow] (P2) -- node[right, font=\footnotesize] {arm별 처치 효과} (OUT);
  \end{tikzpicture}
  \end{center}

  \vspace{-0.4em}
  \footnotesize
  \textbf{Phase 1 중립화의 역할:} (1) 랜덤 배정 검증, (2) 경매 숙련도 균등화 (Plott--Zeiler 2005),
  (3) Algorithm Aversion 차단 (AI 오류 노출 회피, Dietvorst 2015).
\end{frame}

% --- Slide 10 ---------------------------------------------------------------
\begin{frame}{실험 설계 (2) --- 5-arm 처치 구조}
  \begin{center}
  \small
  \begin{tabular}{l|cc}
    & \textbf{정보량 低 (요약)} & \textbf{정보량 高 (상세)} \\
    \hline
    \textbf{Generic} & T1 ($n{=}50$) 비개인화 요약 & T2 ($n{=}50$) 비개인화 상세 \\
    \textbf{AI}      & T3 ($n{=}50$) AI 요약·개인화 & \textbf{T4 ($n{=}50$) AI 상세·개인화} \\
  \end{tabular}
  \end{center}
  \vspace{0.5em}
  $+$ T0 ($n{=}50$): 기본 설명만 (100자) --- 기준선. \textbf{총 $n = 250$.}

  \vspace{0.7em}
  \textbf{식별 핵심 비교:}
  \begin{itemize}\itemsep0.1em
    \item T1 vs T3: AI 개인화 효과 (정보량 통제)
    \item T1 vs T2: 정보량 효과 (비AI 통제)
    \item \textbf{T3 vs T4: 앵커링/과부하 임계점 (AI 통제)} --- 핵심
    \item T4 vs T2: AI vs Generic (고정보량에서)
  \end{itemize}
\end{frame}

% --- Slide 11 ---------------------------------------------------------------
\begin{frame}{실험 설계 (3) --- 제품: 기능성 건강식품}
  \textbf{선정: 프로바이오틱스 1개월분} (또는 오메가-3 / 멀티비타민 후보 검토)

  \textbf{선정 근거:}
  \begin{itemize}
    \item \textbf{신뢰재 (credence good):} 효능 직접 검증 어려움 $\rightarrow$ AI 정보의 불확실성 감소 효과 전제 충족 (채널 A 작동 조건).
    \item \textbf{AI 개인화 자연스러움:} 건강 상태·식이·부족 영양소 $\rightarrow$ 실제 서비스 존재 (Algorithm Appreciation 발생 조건).
    \item \textbf{시장 가격 존재:} 약 ₩30,000/월 $\rightarrow$ 실험 경매 가격대 자연스러움.
    \item \textbf{이중 가상 문제 없음:} 배양육 등과 달리 국내 규제 이슈 없음.
  \end{itemize}

  \vspace{0.3em}
  \muted{단, 순수 credence가 아닌 ``credence-experience 혼합재'' (소화 개선은 1--2주 체감).\\
  $\rightarrow$ 구매 경험 빈도를 채널 A 강도의 \emph{moderator}로 격상 (탐색적).}
\end{frame}

% --- Slide 12 ---------------------------------------------------------------
\begin{frame}{실험 설계 (4) --- AI 메시지 형식}
  \footnotesize
  \begin{tabular}{l|p{4cm}|p{5cm}|p{1.6cm}}
    \hline
    \textbf{Arm} & \textbf{개인화} & \textbf{핵심 내용} & \textbf{Anchor} \\
    \hline
    T0 & 없음 & 기본 설명 100자 & 없음 \\
    T1 (Gen$\times$Lo) & 없음 & 효능 요약 3문장 + 식약처 인증 & 없음 \\
    T2 (Gen$\times$Hi) & 없음 & 성분표 + 임상 근거 + 주의사항 (1쪽) & \textbf{시장 평균가 명시} \\
    T3 (AI$\times$Lo) & 건강 프로필 참조 & 적합도 \% + 권장량 \% (3문장) & \textbf{AI 추정가 명시} \\
    \textbf{T4 (AI$\times$Hi)} & 건강 프로필 + 점수 시각화 & T3 + 성분별 적합도 + 비교 순위 + 발현 시기 + 반응률 (1쪽) & \textbf{AI 추정가 명시} \\
    \hline
  \end{tabular}

  \vspace{0.6em}
  \small
  \key{AnchorDev 직접 측정:}
  \begin{itemize}
    \item T3/T4: $\mathrm{AnchorDev} = \mathrm{Bid} - \mathrm{AI\_price\_estimate}$ (피험자별 $\pm 20\%$ 변이).
    \item T2: $\mathrm{AnchorDev}_{\text{gen}} = \mathrm{Bid} - \mathrm{Generic\_price\_ref}$.
    \item AI 고유 앵커링 $=$ $\mathrm{AnchorDev}_{\text{T4}} - \mathrm{AnchorDev}_{\text{T2}}$.
  \end{itemize}
\end{frame}

% --- Slide 13 ---------------------------------------------------------------
\begin{frame}{실험 설계 (5) --- 경매 메커니즘}
  \textbf{Random $n$-th Price Auction} (Shogren et al.\ 2001):
  \begin{itemize}
    \item 모든 입찰을 순위화 $\rightarrow$ 무작위로 $n$ 추출 $\rightarrow$ 상위 $(n-1)$명이 $n$번째 가격에 구매.
    \item \emph{Off-margin 입찰자}도 engage 시킴 (SPA의 약점 회피).
    \item Phase 1에서 충분한 훈련 후 Phase 2 진입 (Plott--Zeiler 2005 등가 절차).
  \end{itemize}

  \vspace{0.5em}
  \textbf{Filler Task (5분, CVM--경매 사이):}
  \begin{itemize}
    \item CVM 응답이 경매의 \emph{내적 앵커}로 작동하는 순서 효과 차단.
    \item 무관한 인구통계 보충 문항 또는 주의 전환 과제 (인지 부하 없음).
    \item 사후 설문: ``CVM 금액을 경매에서 참고했나'' (1--7점) $\rightarrow$ 이중 앵커링 직접 점검.
  \end{itemize}
\end{frame}

% --- Slide 14 ---------------------------------------------------------------
\begin{frame}{실험 설계 (6) --- 절차와 표본}
  \footnotesize
  \begin{enumerate}\itemsep0.15em
    \item \textbf{사전 설문 (10분):} 건강 프로필, CRT, 친숙도, AI 사용 경험, \textbf{Pre\_Aversion} (Dietvorst 기반).
    \item \textbf{Phase 1 --- Induced (25분):} 메커니즘 교육 + 연습 3R + Induced 경매 5R.\\
      \quad 측정: $\mathrm{Deviation} = |\mathrm{Bid} - \mathrm{TrueValue}|$.
    \item \textbf{Phase 2 --- Homegrown (30분):}
      \begin{enumerate}\itemsep0.05em
        \item Step 1: 처치 자료 (T0--T4)
        \item Step 2: CVM WTP 유도 (가상 구매 시나리오)
        \item Step 2.5: Filler task (5분)
        \item Step 3: 실험경매 WTP (Random $n$-th price, 5R)
        \item Step 4: 사후 설문 (10분) --- Manipulation check, IPV check
      \end{enumerate}
  \end{enumerate}

  \vspace{0.3em}
  \small
  \textbf{표본:} 집단당 $n = 50$, 총 $n = 250$.\\
  \textbf{Effect size 기준:} Fox CVM-X 갭 + Choi et al.\ (2018) 정보 효과 $\pm 15\%$ base rate.\\
  \textbf{소요 시간:} 약 75--80분.
\end{frame}

% =============================================================================
% Part IV: Analysis Strategy (4 slides)
% =============================================================================

% --- Slide 15 ---------------------------------------------------------------
\begin{frame}{분석 전략 (1) --- 주 분석: $2\times2$ 팩토리얼}
  \textbf{Phase 1 (벤치마크 점검):}
  \[
    \mathrm{Deviation}_i = \alpha + \beta_1 T1_i + \cdots + \beta_4 T4_i + \gamma X_i + \varepsilon_i
  \]
  예측: $\beta_1 = \cdots = \beta_4 \approx 0$ (중립화 성공 시).

  \vspace{0.4em}
  \textbf{Phase 2 (주 분석):}
  \begin{align*}
    \mathrm{Bias}_i &= \alpha + \delta_{\text{src}} \cdot \mathrm{AI}_i + \delta_{\text{amt}} \cdot \mathrm{High}_i \\
                    &\quad + \delta_{\text{int}} \cdot (\mathrm{AI}_i \times \mathrm{High}_i)
                              + \gamma \cdot \mathrm{Pre\_Aversion}_i + \delta X_i + \varepsilon_i
  \end{align*}

  \textbf{핵심 가설:}
  \begin{itemize}
    \item $\delta_{\text{src}} < 0$: AI 개인화가 갭 감소 (채널 A)
    \item $\delta_{\text{amt}} < 0$: 정보량 증가가 갭 감소 (채널 A)
    \item \textbf{$\delta_{\text{int}} > 0$: AI $\times$ High = 채널 B 지배 (앵커링/과부하/편향)}
  \end{itemize}
\end{frame}

% --- Slide 16 ---------------------------------------------------------------
\begin{frame}{분석 전략 (2) --- SUR 병렬 (채널 A 비대칭 이동)}
  \textbf{문제:} 채널 A가 작동해도 두 WTP가 \emph{동일 크기로} 이동하면 $\mathrm{Bias}$ 회귀는 효과를 놓침.

  \vspace{0.5em}
  \textbf{Seemingly Unrelated Regression (SUR):}
  \begin{align*}
    \mathrm{WTP}_{\text{CVM},i} &= \alpha_1 + \beta_{1,\text{CVM}} T_i + \gamma_1 X_i + \varepsilon_{1,i} \\
    \mathrm{WTP}_{\text{auc},i} &= \alpha_2 + \beta_{1,\text{auc}} T_i + \gamma_2 X_i + \varepsilon_{2,i}
  \end{align*}

  \textbf{검정:} $H_0\colon \beta_{\text{CVM}} = \beta_{\text{auc}}$ vs $H_1\colon \beta_{\text{auc}} > \beta_{\text{CVM}}$.

  \vspace{0.5em}
  \textbf{해석:}
  \begin{itemize}\itemsep0.1em
    \item $\beta_{\text{auc}} > \beta_{\text{CVM}}$: 경매가 더 크게 반응 $\rightarrow$ 채널 A 확인.
    \item $\beta_{\text{auc}} \approx \beta_{\text{CVM}}$: 두 WTP 동일 이동 $\rightarrow$ 이중 앵커링 가능성.
    \item $\beta_{\text{auc}} < \beta_{\text{CVM}}$: AI가 CVM을 더 올림 $\rightarrow$ 채널 B의 CVM 과장 효과.
  \end{itemize}

  \muted{추가 표본 불필요. 주 분석과 병렬 보고.}
\end{frame}

% --- Slide 17 ---------------------------------------------------------------
\begin{frame}{분석 전략 (3) --- 메커니즘 매개 분석}
  \textbf{경쟁 매개 모형 (Competing Mediation):}\\
  채널 A와 B-1이 관측적으로 등가 $\rightarrow$ 두 매개변수 \emph{동시 투입}:
  \begin{align*}
    \mathrm{WTP}_{\text{auc},i} &= \alpha + \gamma \cdot \mathrm{Treatment}_i \\
                                 &\quad + a_A \cdot \mathrm{Uncertainty\_dec}_i \quad\text{(채널 A)} \\
                                 &\quad + a_B \cdot \mathrm{Anchoring\_score}_i \quad\text{(채널 B-1)} + \delta X_i + \varepsilon_i
  \end{align*}

  \vspace{0.3em}
  \textbf{4 경로 검정:}
  \begin{itemize}\itemsep0.1em
    \item \textbf{경로 A:} Treatment $\rightarrow$ 불확실성 $\downarrow$ $\rightarrow$ $\mathrm{WTP}_{\text{auc}} \uparrow$ $\rightarrow$ Bias $\downarrow$.
    \item \textbf{경로 B-1:} T4 $\rightarrow$ 앵커링 점수 $\uparrow$ $\rightarrow$ $\mathrm{AnchorDev} \approx 0$ $\rightarrow$ Bias 왜곡.
    \item \textbf{경로 B-2:} T4 $\rightarrow$ 과부하 $\uparrow$ $\rightarrow$ 입찰 분산 $\uparrow$ $\rightarrow$ Bias $\uparrow$.
    \item \textbf{경로 B-3:} T3/T4 $\times$ low Pre\_Aversion $\rightarrow$ Appreciation $\uparrow$ $\rightarrow$ 입찰 왜곡.
  \end{itemize}

  \vspace{0.3em}
  \key{Reframing:} ``채널의 완전 분리'' $\rightarrow$ ``두 채널의 \emph{상대적 크기 추정 및 조건부 우세 판별}.''
\end{frame}

% --- Slide 18 ---------------------------------------------------------------
\begin{frame}{분석 전략 (4) --- 이질성과 다중비교}
  \textbf{이질성: $2 \times 2 \times \mathrm{CRT}$ (탐색적, 사전등록)}
  \begin{itemize}
    \item Bergman et al.\ (2010): CRT $\times$ 앵커링 비유의 (p = 0.526). CAT은 유의 (p = 0.043).
    \item $\rightarrow$ CRT 이질성은 \emph{confirmatory 아님}, 탐색적 발견으로 보고.
    \item \textbf{$\varphi_3 \approx 0$이어도 ``Bergman과 일관''이 아닌 ``미탐지(inconclusive)''로 표현.}
  \end{itemize}

  \vspace{0.4em}
  \textbf{구매 경험 Moderator (탐색적):}
  $\psi \cdot (\mathrm{AI}_i \times \mathrm{Experience\_high}_i)$
  $\rightarrow$ 경험자에게 AI 효과 약하면 credence 가정 부분 약화.

  \vspace{0.4em}
  \textbf{다중비교 위계 (사전등록):}
  \begin{itemize}\itemsep0.05em
    \item \textbf{1단계 Confirmatory:} $2 \times 2$의 3개 계수 $\rightarrow$ Holm correction.
    \item \textbf{2단계 Secondary:} SUR + 4경로 mediation $\rightarrow$ FDR (Benjamini--Hochberg).
    \item \textbf{3단계 Exploratory:} $\varphi_3$ (CRT 3-way), $\psi$ (Experience) $\rightarrow$ 보정 없음, 탐색적 보고.
  \end{itemize}
\end{frame}

% =============================================================================
% Part V: Contributions & Discussion (2 slides)
% =============================================================================

% --- Slide 19 ---------------------------------------------------------------
\begin{frame}{예상 기여와 한계}
  \textbf{기여 (6중 진술):}
  \begin{enumerate}\itemsep0.1em
    \item AI 정보의 양면성 \emph{최초 실증} (채널 A vs B 3개 하위 채널).
    \item $2 \times 2$ 팩토리얼로 정보원 $\times$ 정보량 효과 \emph{분리}.
    \item AI 가격 앵커의 WTP 효과 \emph{직접 정량화} (가상편의 맥락 최초).
    \item Algorithm Appreciation $\rightarrow$ WTP 왜곡 \emph{이론화} (advice-taking $\rightarrow$ value elicitation 프레임 설정).
    \item Fox et al.\ (1998) CVM-X 라인의 \emph{21세기 AI 재해석}.
    \item CRT $\times$ 처치 $2 \times 2 \times 2$ 탐색적 이질성.
  \end{enumerate}

  \vspace{0.4em}
  \textbf{한계 (사전 인식):}
  \begin{itemize}\itemsep0.05em
    \item 실험실 $\rightarrow$ 현장 외적 타당성 약함.
    \item AI 금전 앵커 명시가 ecological validity 약화.
    \item 채널 A와 B-1 \emph{관측적 등가성} $\rightarrow$ 완전 분리 불가 (mediation으로 상대 크기만 추정).
    \item $\varphi_3$ underpowered (Simonsohn 2014: interaction $\approx$ $16 \times$ main effect 표본).
  \end{itemize}
\end{frame}

% --- Slide 20 ---------------------------------------------------------------
\begin{frame}{일정과 토론}
  \textbf{일정 (잠정):}
  \begin{itemize}\itemsep0.1em
    \item M+0--1: AI 메시지 prototype $v2$ 작성 + 식약처 가이드라인 검토.
    \item M+1--2: IRB 신청 + 사전등록 (OSF / AsPredicted).
    \item M+2--3: Pilot ($n \approx 30$) $\rightarrow$ 메시지·Anchor 캘리브레이션.
    \item M+3--5: Main study 실행 ($n = 250$).
    \item M+5--7: 데이터 분석 + 논문 작성.
  \end{itemize}

  \vspace{0.5em}
  \textbf{후속 연구 (Future Work):}
  Induced value setting에서 AI source effect를 직접 측정하는 \emph{companion study}\\
  --- Lee et al.\ (2020) 직접 후속 (별도 paper).

  \vspace{0.8em}
  \begin{center}
    \Large \textbf{Discussion}
  \end{center}
\end{frame}

\end{document}
